\section{Analytic comparison with the published total}
\label{sec:compton-analytic-equivalence}

Agreement at finitely many energies is a useful numerical check, but does
not prove that two analytic functions coincide.  This section compares the complete independent result in
Equation~\eqref{eq:comp-total-energy-collected} with Equation (6) of Lee, Schwartz and
Zhang~\cite{qed-processes-lsz}.  The published coefficients enter only this
comparison and are not used to construct the independent amplitude integral.

\paragraph{Observable and normalization.}
Write $s=s_{\mathrm{phys}}/m^2$ and $X=s-1$.  The comparison concerns
unpolarized $e^-\gamma\to e^-\gamma$ together with
$e^-\gamma\to e^-\gamma\gamma$, integrated over the complete photon phase
space.  It excludes the separate charged-pair final state.  Both
calculations use the on-shell mass and physical low-energy charge.  The
published calculation uses dimensional regularization, whereas the
independent construction combines photon-mass and auxiliary Mellin
regulators before removing them.  The comparison is between their finite
sums, not their individual real or virtual finite parts.
The precise normalization is
\begin{equation}
 \sigma_1=\frac{\alpha_{\mathrm T}^3}{m^2}C_{\mathrm{total}}(s),\qquad
 C_{\mathrm{LSZ}}(s)=\frac{\mathcal H_{\mathrm{LSZ}}(X)}{X^3},\qquad X=s-1.
\end{equation}
Here $\mathcal H_{\mathrm{LSZ}}(X)=\mathcal H(X)$ is the numerator defined in the preceding section.  It is
not the energy variable denoted by $H$ in the independent integration.

\paragraph{Classical functions in one iterated-integral basis.}
Let $G$ have the kernels defined in the preceding reference section.
Several changes of basis remove apparent differences immediately:
\begin{align}
 \ln X&=G(0;X),\qquad \ln(1+X)=G(-1;X),\\
 \operatorname{Li}_2(-X)&=-G(0,-1;X),\\
 \operatorname{Li}_2(1-X)-\frac{\pi^2}{6}&=-G(1,0;X),\\
 \operatorname{Li}_3(X^2)-\ln X\operatorname{Li}_2(X^2)
 &=2G(0,1,0;X)+2G(0,-1,0;X).
 \label{eq:compton-equivalence-classical}
\end{align}
For the last identity, apply
$\operatorname{Li}_n(X^2)=2^{n-1}[\operatorname{Li}_n(X)+
\operatorname{Li}_n(-X)]$ and the shuffle product
\begin{equation}
 G(0;X)G(0,a;X)=2G(0,0,a;X)+G(0,a,0;X).
\end{equation}
Alternatively, differentiating both sides of
\eqref{eq:compton-equivalence-classical} gives the same expression, and
both sides vanish at $X=0$.  These arguments initially apply for
$0<X<1$ and extend continuously through $X=1$.  Indeed, the inner kernel
of $G(1,0;X)$ is $\ln X/(X-1)$, which is finite at one.  Thus the
right-hand side is real on the whole positive axis.  On the left-hand
side the imaginary parts cancel when both polylogarithms are continued
on the same side of their cut.  No choice of an imaginary part based on
numerical agreement is required.

After reducing products by shuffle identities, the complete reference
is a linear combination of words of weight at most three with letters
\begin{equation}
 \{0,1,-1,-2,-4,r\},\qquad
 \omega_{r}=\frac{\mathrm{d}X}{\sqrt{X(X+4)}}.
\end{equation}
The letter $1$ is needed for the classical combinations above, although
it is absent from the thirteen nonclassical words displayed separately
in the published formula.

\paragraph{Comparison in the constant alphabet.}
Section~\ref{sec:comp-virtual-energy-basis} derives the constant-letter
representation of the complete virtual contribution from its own
angular primitives. Applying the kernel map
\eqref{eq:compton-equivalence-pullback} to the published total gives
exactly the same set of 116 words. Both expressions use basepoint
$q=1$ and the logarithmic convention fixed there.

The comparison is performed on the complete infrared-finite sums.
No difference between the published total and a partial contribution is
used to reconstruct the remaining sectors.

\paragraph{Equality of the complete finite sum.}
The finite soft reduction in Section~\ref{sec:comp-ck-energy-formula},
the integrated $UW$ term in Equation~\eqref{eq:comp-uw-integrated-compact},
and the other real-emission sectors complete the independent result.
Their sum is Equation~\eqref{eq:comp-total-energy-collected}.
In particular, the intermediate radical $\sqrt{s}$ no longer appears
in its kernels or coefficients. The only square root in the energy
representation is $Y=\sqrt{X/(X+4)}$, with its positive value for $X>0$.
This cancellation follows from the sector reductions before the
published expression is compared.

Let $\mathcal W_{23}$ denote the twenty-three words listed in
Equation~\eqref{eq:comp-total-energy-collected}, including the empty word.
Expanding the published expression with the classical identities above
and the shuffle rule gives exactly the same word set. Thus
\begin{equation}
 C_1(X)-C_{\mathrm{LSZ}}(1+X)
 =\sum_{w\in\mathcal W_{23}}\Delta_w(X,Y)G(w;X),\qquad
 \Delta_w=c_w^{\mathrm{ind}}-c_w^{\mathrm{LSZ}}.
\end{equation}
All twenty-three differences vanish as rational functions of $X$ and $Y$.
No numerical value of $X$, expansion around an endpoint, or relation
between distinct GPL words is used for this final subtraction.
For example, the common coefficients of three words are
\begin{align}
 c_{(1,0)}&=-\frac{X^2-X+2}{X^3},\\
 c_{(0,-1,0)}=c_{(0,1,0)}&=\frac{2(X^2+2X-6)}{X^3},\\
 c_{(r,r,-1)}&=\frac{3(3X^2-8)}{X^3}.
\end{align}
The coefficient comparison includes the empty word, so a residual
additive constant is also excluded. Linear independence of the words
is unnecessary: equality of each coefficient is sufficient.

It remains to identify their definitions on the physical axis.
The root weight in Equations (7)--(8) of the reference is
$Y\,\mathrm dX/X=\mathrm dX/\sqrt{X(X+4)}$, exactly the kernel
$\omega_{r}$ above. Both definitions start at zero, use
$G(0;X)=\ln X$, and integrate along positive real $X$.
All other nonzero letters are negative, except for the two classical
combinations containing the letter $1$.
For those combinations $\ln X/(X-1)$ has a removable singularity.
Their continuous real values were established above by differentiation
and the zero-endpoint boundary values.
The apparently singular elementary term at $X=1$ is likewise fixed by
$\lim_{X\to1}\ln X/(X-1)=1$.
Consequently the coefficient identity gives
\begin{equation}
 \boxed{C_1(X)=C_{\mathrm{LSZ}}(1+X),\qquad X>0,}
\end{equation}
including $X=1$ by continuity. No integration constant or physical
branch has been adjusted to enforce this equality.

The comparison can be reproduced by
\texttt{compton\_total\_reference\_identity.py}.
It reads the independent coefficient dictionary and separately expands
the transcription of Equation (6) of the reference.
The report records every coefficient difference and hashes of both
inputs. The transcription was checked term by term against the published
equation, including its overall $X^{-3}$ normalization and root weight.
The independent calculation does not import that transcription.

\paragraph{A shorter product representation.}
Equality with the published formula does not require keeping its
organization. Section~\ref{sec:comp-total-energy-product} applies further
shuffle identities to the independently obtained coefficients.
Three weight-three words are replaced by products involving lower-weight
functions, and two equal-coefficient words are combined into
$\mathcal P(X)$. The result has an elementary logarithmic polynomial
and thirteen remaining grouped functions. Expanding its products
recovers all twenty-three coefficients exactly.
This reduction decreases the number of separate special-function terms,
although some coefficients then contain elementary logarithms.
It makes no claim that thirteen is a minimal function count.
